\cotefile{14}
\scanpage{101}\mspage{\handwritten{95}}
% archive sheet number: 100
\diagram{\begin{tikzcd}[ampersand replacement=\&]
 E'_{j,n}\handwritten{\arrow[r,"{\handwritten{v_i}}"]} \& E^{\deleted{\unsure{'}}}_{i,n}\\
 E'_{k,n}\handwritten{\arrow[u]\arrow[r,no head,"{\handwritten{v_j}}"]} \& E_{j,n}\handwritten{\arrow[u]}
\end{tikzcd}}
où les flêches verticales sont les homomorphismes de transition.

N.B. On pose $E_{i,n}=E_i/I_i^{n+1}$, $E'_{i,n}=\operatorname{Symm}_{\unsure{C|O}_i}(I_i/I_i^2)/$puissance $(N+1)$ème de l'\unsure{idéah|idéal} d'augmentation. La démonstration du lemme B est immédiate, en procédant comme dans le cas particulier 2.2. déjà traité.

\uline{Lemme} C \unsure{Supposons} $C$ et $B$ \unsure{form} simple sur $A$, (donc en vertu de lemme A les conditions de Lemme B sont vérifiées). Alors:

\handwritten{\unsure{a)}} Pour tout \deleted{\unsure{$i$, existe un $j>i$ et un $v_i$ comme dans lemme B, a)}, avec \uline{en plus} un} système $i,j$, $v_i$ comme dans Lemme B,\inserted{ a)} \deleted{il existe un}\inserted{\handwritten{il existe}} $k\handwritten{\geq}j$ \inserted{\handwritten{et}}\deleted{\unsure{tel que} rendant valable la conclusion de lemme B, b), et tel qu'il existe} un homomorphisme de $A$-algèbres $w_j:E_{k,n}\handwritten{\to}E'_{j,n}$ ayant les propriétés suivantes: 1º) $\handwritten{w_k}$ induit un épimorphisme sur les gradués associés \quad 2º) $v_iw_j$ est l'homomorphisme de transition $E_k\handwritten{\to}E_{\unsure{i}}$

\deleted{b) Soient $i,j,k$, $v_i,v_j$, $w_j$ comme dans a). Alors il existe un $\handwritten{\unsure{l}}\quad k$, un homomorphisme $v_k:E'_{\handwritten{\unsure{l}},n}\quad E_{k,n}$ et un homomorphisme $w_k:E_{\handwritten{\unsure{l}},n}\ \unsure{E}\ E'_{k,n}$ ayant les propriétés suivantes: 1º) $v_k$ a les propriétés mentionnées déjà pour $v_i,v_j$ \quad 2º) $w_k$ a les deux propriétés mentionnées déjà pour $w_j$ \quad \unsure{3}º}
\scanpage{102}
% archive sheet number: 101
\unsure{conjectures|Conjectures} de Lefschetz \unsure{pour} cycles algébriques:
\[
\left\{\begin{aligned}
&\text{Extension aux \uline{Jac. inter.}}\\
&\text{Conjectures \uline{locales}}
\end{aligned}\right.
\]
\scanpage{103}
% archive sheet number: 102
$X$ surface projective non sing.

$Y$ section hyperplane non sing.

\[
\begin{aligned}
\deleted{\unsure{A|D}}\unsure{N|M}^{1}(X)&\longrightarrow\deleted{\unsure{A|D}}N^{1}(Y)
  &&\text{surjectif mod gpe fini}\\
&&&\text{noyau de rang $\rho-1$}\\
J^1(X)&\longrightarrow J^1(Y)&&\text{injectif exc. $p$}\\
A^1(X)&\longrightarrow A^1(Y)&&\text{noyau $K$ tel que $K\cap J^1(X)$ fini}\\
&&&\text{\uline{donc} $K$ de type fini}
\end{aligned}
\]

Soit $B^1(X)$. ``partie primitive'' de $\unsure{A|R|\mathcal{A}}^1(X)$, \unsure{formé|formée} des classes de div. de \uline{degré} $0$, donc
\[
 B^1(X)\longrightarrow J^1(Y)\qquad [B^1(X)\ \text{contient }K]
\]

\diagram{\begin{tikzcd}[ampersand replacement=\&,column sep=large]
\deleted{\unsure{0}}\arrow[r] \& K\arrow[r,"{\substack{\text{inj.}\\\text{mod gpe}\\\text{fini}}}"'] \&
 N^1(X)\arrow[r] \& \underbrace{J^1(Y)/\operatorname{Im}J^1(X)}_{\substack{\text{partie ``Mobile''}\\\text{de }J^1(Y)}}
\end{tikzcd}}

\deleted{\uline{Question} Un \unsure{élmt|élt} de $N^1(X)$ dont l'image dans $J^1(Y)/J^1(X)$ est nulle, est-il dans $K$ i.e. ds. $\operatorname{Im} K$}

N.B. ``En général'' $N^1(X)\longrightarrow J^1(Y)/J^1(X)$ est injectif mod \unsure{gpes} finis, donc $K$ fini.

Je présume qu'il en est ainsi si les positions limites de $Y$ dans les pinceaux \unsure{plans|passant par} contenant $Y$ sont des courbes \uuline{irréductibles} \ldots
\scanpage{104}\mspage{- 115 -}
% archive sheet number: 103
toute sous-algèbre de type fini de $C\handwritten{[}T\handwritten{]}$ vérifie la condition des chaînes, il en est de même de toute sous-algèbre de type fini de $C_{\underline{r}}\handwritten{[}T\handwritten{]}$, ce qui achève la démonstration.

\uline{Corollaire 1 .- Soient} $Y$ \uline{un préschéma irréductible localement noethérien}, $X$ \uline{un préschéma irréductible}, $f:X\handwritten{\to}Y$ \uline{un morphisme} \deleted{\unsure{de type}} \uline{dominant de type fini. Soit} $\handwritten{\xi}$ (resp. $\handwritten{\eta}$) \uline{le point générique de} $X$ (resp. $Y$) \deleted{\illegible{1}}, \uline{et soit} $e=\dim(f^{-1}(\handwritten{\eta}))=\operatorname{deg.tr.}_{\handwritten{\underline{k}}(\handwritten{\eta})}\handwritten{\underline{k}}(\handwritten{\xi})$ (``dimension de la fibre générique\unsure{/}''). \uline{Pour tout} $x\handwritten{\in}X$, \uline{on a, en posant} $y=f(x)$
\[
 e+\operatorname{codim}_Y\handwritten{\overline{\{y\}}}\handwritten{\geq}
 \operatorname{deg.tr.}_{\handwritten{\underline{k}}(y)}\handwritten{\underline{k}}(x)
 +\operatorname{codim}_X\handwritten{\overline{\{x\}}}.\tag*{(13)}
\]
% slanted
\margin{\handwritten{$\dim\underline O_x\leq\dim\underline O_y+\dim\underline O_x\otimes k(y)-\delta(x)$\\
$\delta(x)=\dim_{\deleted{\unsure{y}}x} f^{-1}f(x)-e$}}
% slanted
\margin{\handwritten{formule hybride, donner plutôt \uline{deux} formules, l'une en termes d'alg. comm., l'autre en termes ``géom.''}}

\uline{En outre, les deux membres sont égaux lorsque} $\underline O_y$ \uline{est un quotient d'anneau local régulier}.
% slanted
\margin{\handwritten{\deleted{\illegible{1}} cond. des ch. \unsure{stricte}}}
\margin{\handwritten{\uuline{Variante} \uuline{au cor} $\left\{\begin{aligned}
 \dim\underline O_x&\leq\dim\underline O_y+\dim\underline O_{\deleted{\unsure{y}}x}\otimes\underline k(y)\\
 &\quad-[\dim_x f^{-1}f(x)-e]
\end{aligned}\right.$}}

Comme la question est locale, on peut supposer $X$ et $Y$ affines\inserted{\handwritten{ et réduits}}, \deleted{\illegible{11} \unsure{l'hypothèse} \illegible{5}} \inserted{soient} $X=\operatorname{Spec}(B)$, $Y=\operatorname{Spec}(A)$, où $B$ est une algèbre intègre de type fini sur $A$; comme $f$ est dominant, $A$ s'identifie à une sous-algèbre de $B$ (§ 1, nº2, cor.4 de la prop.5), donc $\underline O_y$ à une sous-algèbre de $\underline O_x$. Comme $\handwritten{\underline k(\eta)}$ et $\handwritten{\underline k(\xi)}$ sont les corps de fractions de $A$ et $B$ respectivement, donc de \deleted{$\unsure{\underline O_x}$ et} $\underline O_y$ et $\underline O_x$ respectivement, la formule (13) résulte de (12) appliquée aux anneaux $\underline O_y$ et $\underline O_x$.
% slanted
\margin{\handwritten{\deleted{Si $\bar x\subset X$ fermé irréductible, soit $\dim\unsure{Z}/f(x)$} \unsure{rédaction voir ? de suite}}}

On notera d'ailleurs que $\operatorname{deg.tr.}_{\handwritten{\underline k}(y)}\handwritten{\underline k}(x)$ n'est autre que la dimension de $f^{-1}(y)\handwritten{\cap\overline{\{x\}}}$, comme il résulte de la prop.7 du nº2, appliquée à $f^{-1}(y)$ considéré comme préschéma algébrique sur $\handwritten{\underline k}(y)$. En particulier:
\margin{\handwritten{\deleted{\illegible{1}}}}

\uline{Corollaire 2} .- \deleted{\unsure{Sous}} \uline{les hypothèses du cor.1 étant vérifiées, supposons en outre que} $x$ \uline{soit fermé dans} $f^{-1}(y)$ \uline{et que} $\underline O_y$ \uline{soit un quotient d'anneau régulier. Alors}
% slanted
\margin{\handwritten{\deleted{\illegible{1}} cond. des chaînes stricte}}
\[
 \operatorname{codim}_X\handwritten{\overline{\{x\}}}
 =e+\operatorname{codim}_Y\handwritten{\overline{\{y\}}}.\tag*{(14)}
\]

\uline{Corollaire 3 .- Les hypothèses du cor.1 étant vérifiées, on a}
\scanpage{106}
% archive sheet number: 105
\uline{Conjectures ``Lefschetz'' locales}

$S=\operatorname{Spec} A$, \qquad $A$ anneau str. local noeth,
\[
 f\in\unsure{rA|\mathfrak r A|\mathfrak{m}_A},
 \qquad
 \left\{\begin{aligned}X&=S-\text{p\textsuperscript{t} fermé}\\Y&=X\cap V(f)\end{aligned}\right.
\]
$X,Y$ réguliers

$U=X-Y$

\textcircled{1}
\[
\begin{aligned}
 \deleted{\unsure{2}}i=\frac N2\quad&
 N^i_{\unsure{\ell}}(X)\equiv\deleted{\unsure{\text{fini}}}0
 &&\text{(mod gpes finis)}\\
 i\geq\frac{N+1}{2}\quad& N^i(X)\deleted{\unsure{\neq}}\equiv0
 &&\inserted{\text{(mod gpes finis)}}\\
 &J^i(X)=0\\
 &\text{i.e.}\quad A^i(X)\equiv0 &&\text{(mod gpes finis)}
\end{aligned}
\]

\textcircled{2}
\[
\begin{aligned}
 J^i(X)\longrightarrow J^i(Y) &\quad\text{bij. \inserted{exc. $p$} si }i\leq\frac{N-2}{2}\\
 &\quad\text{inj \inserted{\unsure{exc.|excp} $p$} si }i\leq\frac{N-1}{2}\\
 &\quad\text{nul si }i\geq\frac N2
\end{aligned}
\]
\margin{[N.B. Si $i\geq\frac N2$, $J^i(Y)=0$ ! d'ailleurs si $i\geq\frac{N+1}2$ on a $J^i(X)=J^i(Y)=0$]}

\[
\begin{aligned}
 N^i(X)\longrightarrow N^i(Y)&\quad\text{bij \inserted{exc. $p$} si } i\leq\frac{N-3}2\\
 &\quad\text{inj \inserted{exc. $p$} si } i\leq\frac{N-2}2\\
 &\quad\text{nul \inserted{mod. gpes finis} si } i\geq\frac{N-1}2
\end{aligned}
\]
\margin{Cas critique \uuline{$J^{\frac{N-1}2}$} $J^{\frac N2}$}
\margin{[N.B si $i\geq\frac{N-1}2$, on a $N^i(Y)\equiv0$ ! si $i\geq\frac N2$, \unsure{on a} $N^i(X)=N^i(Y)\unsure{\equiv|=}0$ !]}
\scanpage{107}
% archive sheet number: 106
Cas critique \uline{$\deleted{\unsure{J}}N^{\frac{N-2}2}$}, $\deleted{\unsure{J}}N^{\frac{N-1}2}$

\[
\begin{aligned}
 A^i(X)\unsure{\mathord{-}|\to} A^i(Y)&\quad\text{bij si }i\leq\frac{N-3}2\\
 &\quad\text{inj. si }i\leq\frac{N-2}2
\end{aligned}
\]
Cas critiques $i=\dfrac{N-2}2,\dfrac{N-1}{2},\dfrac N2$

% slanted
\margin{conséquence de 1º et 2º}
\textcircled{3}
\[
\begin{aligned}
 A^i(X)\longrightarrow A^i(U)&\quad\text{\uline{un isom} mod gpes finis, }\forall i\\
 &\quad\text{\uline{un} \uline{isom} si } i\leq\frac{N-1}2
\end{aligned}
\]
\margin{[N.B. si $i\geq\frac{N+1}2$, $\deleted{\unsure{i\geq N}}A^i(X)$ et $A^i(U)$ \unsure{sont|est} fini !]}

\[
\begin{aligned}
 A^i(Y)\longrightarrow A^{i+1}(X)&\quad\text{\uline{nul} \unsure{mod} \deleted{\unsure{f}} gpes finis, }\forall i\\
 &\quad\text{nul si } i\leq\frac{N-1}2
\end{aligned}
\]
\margin{[N.B si $i\geq\frac N2$, $A^i(Y)$ est un gpe fini !]}

\textcircled{4} \uline{Pb} Définir des formes ``\unsure{intersections|intersection}'' $N^i(X)\times N^i(X)\longrightarrow\mathbb Z$ ??

Définir pour $J^{\frac{N-1}2}(X)\longrightarrow J^{\frac{N-1}2}(Y)$ resp. $N^{\frac{N-2}2}(X)\longrightarrow N^{\frac{N-1}2}(Y)$ une théorie des cycles évanouissants \enspace - - - - - -

\noindent\rule{\linewidth}{0.35pt}

\uline{Ex}\quad $N=2$

\textcircled{1}\quad \fbox{$N^1(X)$ fini}: c'est le th. \uline{de \deleted{\illegible{1}} Mumford-Hodge}.

\textcircled{2}\quad vide

\textcircled{3}\quad $\deleted{\unsure{A^1(X)|\mathcal{A}^1(X)}\longrightarrow}$ vide

Groupe \deleted{\unsure{qui|fini}} essentiel \qquad $J^1(X)$
\scanpage{108}
% archive sheet number: 107
\uline{Ex}\quad $N=3$

\textcircled{1}\quad \fbox{$A^2(X)$ fini} \quad en particulier $J^{\unsure{2|1}}(X)=0$

\textcircled{2}\quad \fbox{$J^1(X)\longrightarrow J^1(Y)$ \quad injectif exc. $p$}\deleted{\unsure{injectif}}

\textcircled{3}\quad \deleted{$\unsure{A^2(X)|\mathcal{A}^2(X)}$} contenu dans 1º)

groupes essentiels \quad $J^1(X),N^1(X)$ \qquad définissant $A^1(X)$

\unsure{hom|homs} essentiels
\[
\begin{aligned}
 &J^1(X)\longrightarrow J^1(Y)\\
 &N^1(X)'\longrightarrow J^1(Y)
\end{aligned}\ \bigg|\ A^1(X)'\longrightarrow J^1(Y)
\]
% slanted
\margin{sous-gpe d'indice fini de $N^1(X)$}

\uline{Pb} définir un complémentaire naturel de $J^1(X)$ dans $J^1(Y)$ (correspondant aux ``cycles \unsure{évanescents|évanescent}'')

Donner un sens à la \unsure{formule}
\[
 J^1(X)\overset{\sim}{\longrightarrow}J^1(Y)^{\pi}
\]
Interpréter le th. de \inserted{l'index de} Hodge pour les surfaces projectives en termes de la négativité d'une certaine forme définie sur $N^1(X)$\ldots?

\uline{Ex}\quad $N=4$

\textcircled{1}
\[
 \left\{\begin{aligned}
 &N^2(X)\ \text{fini}\\
 &A^3(X)\ \text{fini, en particulier }J^3(X)=0
 \end{aligned}\right.
\]
\textcircled{2}
\[
\begin{aligned}
 J^1(X)&\longrightarrow J^1(Y)&&\text{bij exc. $p$}\\
 \deleted{\unsure{J}}N^1(X)&\longrightarrow N^1(Y)&&\text{inj. exc $p$}\\
 &\Rightarrow\quad A^1(X)\longrightarrow A^1(Y)&&\text{\uline{bij exc.} $p$}
\end{aligned}
\]
\deleted{\textcircled{3}}
\scanpage{109}
% archive sheet number: 108
groupes essentiels \qquad $J^2(X),N^1(X)$

hom essentiel \qquad $N^1(X)\unsure{\mathord{-}|\to}N^1(Y)$

\uline{Pb} définir un complémentaire naturel de $N^1(X)$ dans $N^1(Y)$ (cycles évanouissants). \deleted{Donner}

Donner un sens à la formule
\[
 \deleted{\unsure{J}}N^1(X)\overset{\sim}{\longrightarrow}N^1(Y)^{\pi}
\]

\noindent\rule{\linewidth}{0.35pt}

Cas important \qquad $N=2m+1$, à \unsure{prouver} alors
\[
 A^{m+1}(X)\equiv0\qquad(\text{et }A^i(X)\equiv0\text{ si } i\geq m+1)
\]
[Cas \qquad $N=2m$, alors
\[
\begin{aligned}
 N^m(X)&\equiv0\quad\text{i.e.}\quad J^m(X)\equiv A^m(X),\\
 \text{et}\quad A^{m+1}(X)&\equiv0\ \unsure{\sim|,}\quad
 (\text{et }A^{\deleted{\unsure{m}}i}(X)\equiv0\text{ si }i\geq m+1)
\end{aligned}
\]
\scanpage{111}
% archive sheet number: 110
\uline{Conjectures} \inserted{``\uline{Lefschetz}''} \quad ($k$ sép. clos)

Soit $X^n$ projectif lisse connexe, $Y^{n-1}$ section hyperplane lisse.

\[
\begin{aligned}
 -\quad A^i_r(X)&\ \ldots\quad\text{cycles de codimension $i$ mod. équivalence \deleted{alg} \inserted{rationnelle}}\\
 A^i_a(X)&\hspace{3em}\text{algébrique}\\
 A^i_{\tau}(X)&\hspace{3em}\text{\unsure{équivalence} $\tau$}
\end{aligned}
\]

\textcircled{1}\quad $A^i_r(X)\longrightarrow A^i_r(Y)$
\[
\boxed{\begin{aligned}
 \deleted{\text{bijectif \inserted{exc. $p$} si }\ 2i\leq n-2\quad\text{i.e. }i<\frac{n-1}2}\\
 \deleted{\text{injectif exc. $p$ si }\ 2i\unsure{\leq|=}n-1\quad\text{i.e. }i=\frac{n-1}2}\\
 \deleted{\text{bijectif \unsure{sur les|exc. $p$}} \ \unsure{A^i_r}\ \text{\unsure{si|sur}}}
\end{aligned}}
\]
\[
\begin{aligned}
 2i\leq n-2&\qquad\underset{\uline{\unsure{e|\leq}}}{\text{i.e.}}\  i<\frac{n-1}2\\
 &\text{bijectif [exc. $p$ ?]}\\
 2i=n-1&\qquad\text{i.e. }i=\frac{n-1}2\\
 &\text{injectif [exc. $p$,]}\\
 &\text{induit}\quad A^i_r(X)^\tau\longrightarrow A^i_r(Y)^\tau\quad\text{bijectif [exc. $p$]}\\
 &\Rightarrow\quad A^i_a(X)\longrightarrow A^i_{\unsure{a|r}}(Y)\quad\text{bijectif [exc. $p$]}\\
 &\phantom{\Rightarrow\quad} A^i_\tau(X)\unsure{\mathord{-}|\to}A^i_\tau(Y)\quad\text{injectif}\\
 2i=n\deleted{\unsure{-1}}&\qquad\text{i.e. }i=\frac{n\deleted{\unsure{-1}}}2\\
 &\text{induit}\quad A^i_r(X)^\tau\longrightarrow A^i_{\unsure{r|a|\tau}}(Y)^\tau\quad\text{injectif [exc. $p$]}
\end{aligned}
\]

\textcircled{2}\quad $A^r_i(Y)\longrightarrow A^r_i(X)$
\[
\begin{aligned}
 2i\leq n-3&\qquad\text{i.e. }i\leq\frac{n-3}2\deleted{\unsure{=(n-1)/2+1}}=\frac{n-1}2-1\\
 &\text{bijectif [exc. $p$ ?]}\\
 2i=n-2&\qquad\text{i.e. }i=\frac{n-2}2
\end{aligned}
\]
\[
 \left\{\begin{aligned}
 &\text{surjectif [exc. $p$,]}\\
 &\text{induit}\quad A_i^r(Y)^\tau\longrightarrow A_i^r(X)^\tau\qquad\text{surj. [exc. $p$]}\\
 &\text{induit}\quad A_i^a(Y)\longrightarrow A_i^a(X)\qquad\text{bij. [exc. $p$]}
 \end{aligned}\right.
\]
\scanpage{112}
% archive sheet number: 111
\[
\begin{aligned}
 2i=n-1&\qquad\text{i.e.}\qquad i=\frac{n-1}2\\
 &\text{induit}\qquad A_i^\tau(Y)\unsure{\mathord{-}|\to} A_i^\tau(X)\qquad\text{surj. exc. $p$}
\end{aligned}
\]

\textcircled{3}
\[
\begin{aligned}
 A_a^i(X)&\xrightarrow{E^{n-2i}} A_a^{n-i}(X)\qquad(2i\leq n)\\
 &\text{\deleted{injectif exc. $p$}\inserted{bijectif} mod groupes finis}
\end{aligned}
\]
\deleted{induit $\deleted{\unsure{B}}A_\tau^i(X)\longrightarrow A_\tau^{n-i}(X)$ bijectif mod groupes finis}
\[
\begin{aligned}
 A^i_{\deleted{\unsure{a|\tau}}r}(X)^\tau
 &\xrightarrow{E^{n-2i+1}} A_r^{n-i+1}(X)^\tau\qquad 2i-1\leq n\\
 &\text{bijectif mod groupes finis}
\end{aligned}
\]
\[
 \Longrightarrow\qquad A^i_r(X)\xrightarrow{E^p}A^{i+p}_r(X)
\]
\unsure{a)} injectif mod gpes finis si $p\leq n-2i$

b) surjectif mod gpes finis si $p\geq n-2i+1$

\noindent\rule{\linewidth}{0.35pt}
\scanpage{113}
% archive sheet number: 112
\textcircled{$2'$} Posons $U=X-Y$, qui est une variété affine non sing.
\[
 A^i_r(U)\simeq A^i_r(X)/\operatorname{Im} A^i_r(Y),\quad
 \text{de }=\text{ pour } A^i_a(U)\ldots
\]
\[
\begin{aligned}
 2i\leq n-2&\quad\text{i.e. }i\leq\frac n2-1&
 A_i^r(U)&=0\quad(\text{exc. $p$ ?})\deleted{\unsure{??}}\\
 2i=n-1&\quad\text{i.e. }i=\frac{n-1}2\deleted{\unsure{+1}}&
 A_i^a(U)&=0\quad(\text{exc $p$ ?})
\end{aligned}
\]

\textcircled{$3'$} Soit $E$ le cône projetant affine épointé de $X$, (variété quasi-affine non sing. de dim $2n+1$).

\[
\begin{aligned}
 2i\geq n+2&\quad\text{i.e. } i\geq\frac n2+1&
 A_r^i(E)&\ \deleted{\unsure{=0|\equiv}}\text{ fini},\quad A^i_r(E)^a=0\\
 &\quad n-i\leq\frac n2-1\\
 2i\geq n+1&\quad\text{i.e. } i\geq\frac{n+1}2&
 A^i_a(E)&\ \text{fini}\\
 &\quad\text{i.e. }\unsure{n|N}-i\leq\frac{n-1}2
\end{aligned}
\]

\uline{N.B.} On doit avoir
\[
\left\{\begin{aligned}
 A_a^i(E)&\simeq\deleted{\unsure{B_a^i(X)}}\unsure{N|M|P}_a^i(X)
 \quad\text{\inserted{mod groupes finis} si }2i\leq n\\
 \deleted{A_r^i(E)^a}&\deleted{\simeq\unsure{B_r^i|P^i_{\tau}}(X)\quad\text{mod gpes finis si }2i\leq n+1}\\
 A_a^i(E)&\simeq0\quad\text{mod gpes finis si }2i\geq n+1
\end{aligned}\right.
\qquad
\]
\[
\left\{\begin{aligned}
 A_r^i(E)^a&\simeq I^{i}_{\deleted{\unsure{B_r}}}(X)\quad\text{mod gpes finis si }2i\leq n+1\\
 A_r^i(E)^{\unsure{a|-}}&\simeq\unsure{0|B}\qquad\text{``}\qquad\text{si }2i\geq n+2
\end{aligned}\right.
\]
où
\[
\begin{aligned}
 \deleted{\unsure{B}}N^i_a(X)&\simeq\operatorname{Ker}\bigl(A^i_a(X)\xrightarrow{L^{n-2i}}A^{n-i}_{\unsure{a}}(X)\bigr)\\
 I^i_{\unsure{r|a}}(X)&=\operatorname{Ker}\bigl(A^i_r(X)^a\xrightarrow{L^{n-2i+1}}A^{n-i+1}_r(X)^a
\end{aligned}
\]
\margin{partie ``primitive'' de la ``\unsure{coh.} algébrique''}
\margin{partie ``primitive'' \deleted{de} la V.A. intermédiaire}
\scanpage{114}
% archive sheet number: 113
\deleted{\uline{Th. 4} Soit $A$ un anneau local noeth complet \deleted{\unsure{dont}}\unsure{tel} que les composantes irréductibles de $S=\operatorname{Spec} A$ sont ``d'égales \unsure{car|car.|dim}''. Soit $f\in\unsure{rA|\mathfrak r A|\mathfrak{m}_A}$, et $U=S_f$. Alors \deleted{\unsure{pour tt entier}} $\operatorname{cd} U\leq\dim A$.

On peut supposer $A$ intègre \inserted{normal} complet, \inserted{de dim $n$}, et on est ramené à prouver que $H^i(U,\mathbb Z/\ell\mathbb Z)=0$ si $i>n$. $A$ est fini sur un anneau de \deleted{\unsure{séries}} formelles $A_0=k[[t_1,\ldots,t_n]]$, et on peut supposer d'abord $K/K_0$ \uline{séparable} (quitte à remplacer $A$ par la normalisée de $A_0$ dans la clôture sép. de $K_0$ dans $K$, \ldots). \inserted{Soit $\delta\in rA_0$ tel que $S_{\delta,0}$ \unsure{étale sur} $S_{0,\delta}$.} On \deleted{\unsure{peut écrire le th.} \illegible{1}} peut supposer $\delta f=t_n$ [par la dém. du th. de Cohen \ldots].

\diagram{\begin{tikzcd}[ampersand replacement=\&,row sep=small]
 S\arrow[d,no head] \&\\
 S_0=\operatorname{Spec} k[[t_1,-t_n]]\arrow[dr] \&\\
 \& S'_0=\operatorname{Spec} k[[t_1-t_{n-1}]]\{t_n\}
\end{tikzcd}}

\deleted{\unsure{Posant|Posons}} $A'_0=\deleted{\unsure{\operatorname{Spec}}} k[[t_1,-t_{n-1}]]\{t_n\}$, $S'_0=\operatorname{Spec} A'_0$, le th. 3. \deleted{\illegible{1}} montre que $S\mid\unsure{D_0|S_0}_{\delta f}$ étale provient par changt de base d'un $Z'$ étale sur $S'_{0\,t_n}$ et en fait (par raison de \unsure{normalité|connexité})
\[
 S\simeq S_0\times_{S'_0}S',\qquad S'\inserted{=\operatorname{Spec} A'}\ \text{rev. fini local de }S_0.
\]
D'ailleurs $f$ provient d'un $f'\in rA'$.}
\scanpage{115}
% archive sheet number: 114
\deleted{1) \unsure{$\dim$} \unsure{$X$} $n=2$}

\uline{Ex}, \qquad $n=2$

\textcircled{1}
\[
\begin{aligned}
 A_r^0(X)&\mathord{-}A_r^0(Y)&&\text{bijectif exc $p$}&&\text{Ok}\\
 A_r^1(X)^\tau&\longrightarrow A_r^1(Y)^\tau&&\text{injectif exc. $p$}&&\text{Ok}
\end{aligned}
\]
\textcircled{2}
\[
\begin{aligned}
 A^r_0(Y)&\longrightarrow A^r_0(X)&&\text{surjectif exc. $p$}&&\text{ok}\\
 A^r_0(Y)^\tau&\longrightarrow A^r_0(X)^\tau&&\text{- - - ---}&&\text{Ok}\\
 \underset{\substack{\Vert\\\mathbb Z}}{A^a_0(Y)}&\longrightarrow
 \underset{\substack{\Vert\\\mathbb Z}}{A^a_0(X)}&&\text{bijectif exc. $p$}&&\text{Ok}
\end{aligned}
\]
\margin{du moins si équiv. rat. des 0-cycles $=$ équiv. d'Albanese\ldots}

\textcircled{3}
\[
\left.\begin{aligned}
 \underset{\substack{\Vert\\\mathbb Z}}{A_a^0(X)}&\xrightarrow{E^2}
 \underset{\substack{\Vert\\\mathbb Z}}{A_a^2(X)}\\
 A_a^1(X)&\overset{E^0=\mathrm{id}}{\mathord{-}}A_a^1(X)
\end{aligned}\right\}\quad\text{bijectif mod gpes finis}\quad\text{O.K.}
\]
\[
\left.\begin{aligned}
 \underset{\substack{\Vert\\0}}{A_r^0(X)^\tau}&\longrightarrow
 \underset{\substack{\Vert\\0}}{A_{\unsure{a|r}}^2(X)}\\
 A_r^1(X)^\tau&\xrightarrow{\mathrm{id}}A_r^1(X)^\tau
\end{aligned}\right\}\quad\text{bij. mod gpes finis}\quad\text{O.K.}
\]

\textcircled{$2'$}\qquad $A_0^r(U)=0$
\margin{O.K. si équivalence rationnelle $=$ équiv. d'Albanese}

\uline{Ex}.\quad $n=3$

\textcircled{1}
\[
\begin{aligned}
 A_r^0(X)&\mathord{-}A_r^0(Y)&&\text{\unsure{bij.|bij} exc $p$}&&\text{O.K.}\\
 A_r^1(X)&\longrightarrow A_r^1(Y)&&\text{injectif exc $p$.}&&\text{O.K.}\\
 A_r^1(X)^\tau&\mathord{-} A_r^1(Y)^{\unsure{\tau|2}}&&\text{bijectif exc $p$}&&\text{O.K.}
\end{aligned}
\]
{\scriptsize
\diagram{\begin{tikzcd}[ampersand replacement=\&,column sep=3pt,row sep=small]
 \text{\textcircled{2}}\quad A_0^r(Y)\longrightarrow A_0^r(X)
 \& \text{bijectif exc. }p \& \text{O.K.}
 \& \substack{\text{si équiv. rat.}\\=\text{équiv. Alb.}}\\
 A_1^a(Y)\longrightarrow A_1^a(X) \& \text{surjectif exc. }p \& \text{\textcircled{?}} \&\\
 \text{\textcircled{3}}\quad \underset{\substack{\Vert\\\mathbb Z}}{A^0_a(X)}
 \overset{E^3}{\mathord{-}}\underset{\substack{\Vert\\\mathbb Z}}{A^3_a(X)}
 \& \text{bijectif mod gpes finis}\arrow[rr,bend right=12,no head] \& \text{O.K.}
 \& \substack{\text{\unsure{de cod|décid.}}\\\text{g\textsuperscript{pes} finis}}\\
 A^1_a(X)\overset{E^1}{\mathord{-}}A^2_a(X)
 \& \text{\textcircled{?}}\arrow[uur,Rightarrow]\arrow[d,no head]\ \text{\deleted{bijectif}}
 \& \text{bijectif mod gpes finis} \&\\
 \& \text{Question ``numérique''} \& \&
\end{tikzcd}}}
\scanpage{116}
% archive sheet number: 115
\[
\begin{aligned}
 \underset{\substack{\Vert\\0}}{A_r^0(X)^\tau}&\xrightarrow{E^4}
 \underset{\substack{\Vert\\0}}{A_r^{4\deleted{\unsure{+}}}(X)^\tau}
 &&\text{\unsure{bij exc.|bijexe.|bijec.} gpes finis}&&\text{Ok}\\
 A_r^1(X)^\tau&\xrightarrow{E^2}A_r^3(X)^\tau&&\text{id}&&\text{Ok}\\
 A_r^2(X)^\tau&\overset{\mathrm{id}}{\mathord{-}}A_r^2(X)^\tau&&\text{id}&&\text{ok}
\end{aligned}
\]
\margin{si équiv. \unsure{rat.} $\simeq$ équiv. Alb.}

\textcircled{$2'$}\qquad $A_0^r(U)=0\qquad A_1^a(U)=0$ \qquad (exc. $p$) \quad \textcircled{??}

\uline{Ex}.\quad $n=4$

\textcircled{1}
\[
\begin{aligned}
 \underset{\substack{\Vert\\\mathbb Z}}{A_r^0(X)}&\mathord{-}
 \underset{\substack{\Vert\\\mathbb Z}}{A_r^0(Y)}&&\text{bij. exc. $p$}&&\text{ok}\\
 A_r^1(X)&\mathord{-} A_r^1(Y)&&\text{bij. exc. $p$}&&\text{O.K.}\\
 \deleted{\unsure{A^2_r(X)}}&\deleted{\unsure{\longrightarrow A^2_r(Y)}}
\end{aligned}
\]
\diagram{\scriptsize\begin{tikzcd}[ampersand replacement=\&,row sep=small,column sep=2pt]
 {} \& A_r^2(X)^\tau\unsure{\mathord{-}|\to}A_r^2(Y)^\tau \& \text{injectif exc. $p$} \& |[draw,circle,inner sep=1pt]| ? \& \text{question ``continue''}\\
 \text{\textcircled{2}} \& A^r_0(\unsure{Y|X})\mathord{-} A^r_0(X) \& \text{bijectif exc. $p$} \& \text{O.K.} \& \text{si équiv. rat. $=$ équiv. Alb}\\
 {} \& A^r_1(Y)\longrightarrow A^{\unsure{r}}_1(X) \& \text{surj. exc $p$} \& ? \&\\
 {} \& A^r_1(Y)^\tau\longrightarrow A^r_1(X)^\tau \& \text{surj. exc. $p$} \& |[draw,circle,inner sep=1pt]| ? \& \text{question ``continue''}\\
 {} \& A^a_1(Y)\longrightarrow A^a_1(X) \& \text{bij. exc. $p$.} \& |[draw,circle,inner sep=1pt]| ? \& \substack{\text{question ``numérique''}\\\text{(ess.\textsuperscript{t})}}\\
 \text{\textcircled{3}} \& \underset{\substack{\Vert\\\mathbb Z}}{A_a^0(X)}\overset{E^4}{\mathord{-}}\underset{\substack{\Vert\\\mathbb Z}}{A_a^4(X)} \& \text{bij. mod gpes finis} \& \text{O.K.} \& \text{mod gpes finis}\\
 {} \& A_a^1(X)\xrightarrow{E^2}A_a^3(X) \& \text{id} \& |[draw,circle,inner sep=1pt]| ?\arrow[uu,Rightarrow,bend left=18] \&\\
 {} \& A_a^2(X)\overset{\mathrm{id}}{\mathord{-}}A_a^2(X) \& \text{id}\arrow[uurr,no head,bend right=10] \& \text{Ok.} \&\\
 {} \& \underset{\substack{\Vert\\0}}{A_r^0(X)^\tau}\xrightarrow{E^5}\underset{\substack{\Vert\\0}}{A_r^5(X)^\tau} \& \text{\deleted{bij.} id} \& \text{ok} \&\\
 {} \& A_r^1(X)^\tau\xrightarrow{E^3}A_r^4(X)^\tau \& \text{id} \& \text{Ok} \& \text{si équiv. rat. $=$ équiv. Alb}\\
 {} \& A_r^2(X)^\tau\xrightarrow{E}A_r^3(X)^\tau \& \text{id} \& |[draw,circle,inner sep=1pt]| ?\arrow[uuuuuuu,Rightarrow,bend left=23]\arrow[uuuuuuuuuu,Rightarrow,bend left=34] \& \text{question ``continue''}
\end{tikzcd}}

\textcircled{$2'$}\qquad $A_0^r(U)=0,\qquad A_1^r(U)=0$\qquad exc. $p$.\qquad \textcircled{??}
